\section{技术栈收敛与 AI 友好性}

\subsection{从``开发者品味''到``AI 友好性''}

Birgitta 提出了一个大胆的预测：随着编程从``手动编写代码''转向``引导代码生成''，AI 可能推动我们走向\textbf{更少的技术栈}。

她观察到，框架和 SDK 的易用性（usability）仍然重要------反复的实践表明，\textbf{对人类友好的东西对 AI 也友好}。但在具体细节层面，\textbf{开发者的个人品味将变得不那么重要}。接口中的小低效和特殊习惯不再那么烦人，因为我们不再直接与之打交道。

\begin{knowledgebox}{AI 驱动的技术栈选择标准变迁}
在传统开发中，技术栈选择往往受以下因素影响：
\begin{itemize}[leftmargin=1.5em]
    \item 团队熟悉度和开发者偏好
    \item 生态系统成熟度和社区活跃度
    \item 性能特征和部署便利性
    \item API 设计的人体工学（ergonomics）
\end{itemize}

在 AI 驱动的未来中，新的选择标准可能包括：
\begin{itemize}[leftmargin=1.5em]
    \item 是否有高质量的 harness 模板可用
    \item 代码模式是否容易被 AI 理解和生成
    \item 架构约束是否容易通过工具实施
    \item ``AI 友好性''------AI 在该技术栈上的生成质量
\end{itemize}
\end{knowledgebox}

\subsection{代码库拓扑作为新的抽象层}

Birgitta 提出了一个更深层的问题：如果我们能够广泛地 harness 代码库设计模式，那么这些\textbf{拓扑结构}（topologies）是否会成为新的抽象层------而非许多 AI 爱好者所期望的\textbf{自然语言}？

OpenAI 团队讨论了\textbf{架构刚性}（architectural rigidity）和\textbf{强制规则}（enforcement rules），主要关注两个领域：

\begin{itemize}[leftmargin=1.5em]
    \item \textbf{保持数据结构稳定}：确保核心数据模型不被随意修改
    \item \textbf{定义和强制模块边界}：确保模块之间的接口清晰且稳定
\end{itemize}

Birgitta 承认这些听起来合理，但缺乏具体示例让她难以想象 ``we require Codex to parse data shapes at the boundary'' 在实践中具体是什么样子。

\subsection{本章小结}

AI 可能推动技术栈的收敛------团队将优先选择有良好 harness 支持的技术栈。代码库的拓扑结构可能取代自然语言，成为 AI 编程时代的关键抽象层。这一趋势将重塑我们选择和评估技术的方式。

\newpage

